%!TEX root = ../main.tex
% Abstract: traducción fiel del resumen. Las keywords se definen en main.tex.

\begin{abstract}

The care of a person with a disability usually falls on relatives with no
specific training, on hired caregivers or on the staff of an institution.
Caregivers must organize medication, appointments, medical tests and
paperwork, and the information they need is scattered across papers,
messages and official websites that are hard to navigate. When several
people share the care, there is also no common record of what was done and
by whom.

This work presents ``Acompañar'', a web platform adapted to mobile phones that
centralizes care management and organizes the people involved in groups
with an owner, caregivers who record the daily care and observers who can
only consult it. Its core component is a hybrid and explainable recommender
system. It only considers resources for the person's type of disability,
compares their profile with a verified catalog, discards what does not
apply, incorporates the ratings of other caregivers and explains every
suggestion. It also includes a step-by-step guide to the disability
certificate (CUD) based on official regulations, an emergency module and a
subscription model with a free plan and two paid plans. The system was built
with React, ASP.NET Core and SQL Server. It currently runs in a local
environment, and its deployment on Azure is proposed as a next step.

The evaluation included 68 automated tests, an analysis of the
recommendations over thirteen demonstration profiles, with 100\% of the
suggestions explained, the verification of 52 links against official
sources and a survey analyzed over 80 cases generated from a test. The
contribution of this work is the integration, in a single place and adapted
to the Argentine context, of care organization, group coordination and
explainable guidance. Validation with real users is the next step.

\end{abstract}
